\begin{frame}{생성 속도 행: ``안정성의 대가''}
  \begin{itemize}
    \item 가장 \textbf{안정한} Diffusion은 생성에 \textbf{$T$회
      반복} 필요 (DDIM으로 줄여도 수십 회)
    \item 가장 \textbf{빠른} GAN/VAE는 \textbf{한 번의} 순전파로 생성
  \end{itemize}
  \vskip0.8em
  \begin{block}{장 전체를 관통하는 트레이드오프}
    ``\textbf{데이터를 설명할 수 있다}(안정성, 평가 가능)'' vs
    ``\textbf{빨리 만들 수 있다}(속도)''
  \end{block}
  \vskip0.6em
  \begin{block}{실전 시스템은 ``하나씩 고르는 것''이 아니라 \textbf{조합}}
    Stable Diffusion = \kb{VAE}(잠재 공간) $\times$
    \kb{Diffusion}(노이즈 제거). \;
    SceneDiffuser = \kb{Diffusion} 원리를 \textbf{시간 축의 장면}으로
    확장. \; $\Rightarrow$ ``세 원리를 나란히 보고 나서야,
    ``관측되지 않은 구조가 데이터를 만들어낸다''는 하나의 질문이
    얼마나 다양한 방식으로 풀릴 수 있는지''가 보임.
  \end{block}
\end{frame}
